% TOP_PROBLEM: {"id": 30007337, "title": "Full-boundary recovery of both smooth isotropic elastic moduli", "category_id": 9, "category": {"id": 9, "name": "pde", "display_name": "Partial Differential Equations", "description": "PDEs and their applications in physics and geometry.", "slug": "pde", "order_index": 9, "created_at": "2026-07-31T15:26:25.671Z"}, "status": "open"}
% FIRST_POSTED: 2026-10-06
% ENTRY_KIND: statement_only
% !TeX program = lualatex
\documentclass[11pt]{article}
\usepackage[margin=25mm]{geometry}
\usepackage{fontspec}
\setmainfont{Latin Modern Roman}
\usepackage{amsmath,amssymb,mathtools}
\usepackage{unicode-math}
\setmathfont{Latin Modern Math}
\usepackage{xurl,hyperref}
\hypersetup{colorlinks=true,urlcolor=blue}
\setcounter{secnumdepth}{0}
\setlength{\parindent}{0pt}
\setlength{\parskip}{5pt}
\setlength{\emergencystretch}{3em}
\newcommand{\R}{\mathbb R}
\newcommand{\N}{\mathbb N}
\newcommand{\Z}{\mathbb Z}
\newcommand{\Q}{\mathbb Q}
\newcommand{\C}{\mathbb C}
\newcommand{\D}{\mathbb D}
\newcommand{\F}{\mathbb F}
\newcommand{\sref}[2]{\hyperlink{source:#1:#2}{[S#2]}}
\newcommand{\cataloguescope}[1]{#1}
\begin{document}
\section*{Full-boundary recovery of both smooth isotropic elastic moduli}\label{30007337}
\subsection{Definitions and mathematical statement}
\paragraph{Statement recorded in the supplied source.}
For any bounded connected smooth domain in $R^3$, do identical static isotropic-elasticity traction Dirichlet-to-Neumann maps on $H^{1/2}(boundary;R^3)\to H^{-1/2}(boundary;R^3)$ imply equality of both smooth real Lamé functions lambda and mu throughout the domain, assuming $\mu >0$ and $3\lambda +2\mu >0$?
\subsection{Short English statement}
For any bounded connected smooth domain in $R^3$, do identical static isotropic-elasticity traction Dirichlet-to-Neumann maps on $H^{1/2}(boundary;R^3)\to H^{-1/2}(boundary;R^3)$ imply equality of both smooth real Lamé functions lambda and mu throughout the domain, assuming $\mu >0$ and $3\lambda +2\mu >0$?
\subsection{Sources}
\begin{itemize}
\item[S1] Original question source or mathematical context.
\url{https://github.com/openai/math/blob/adc7f1241b42e322a6451854ab7e4b4c146bf78a/preprints/Global-Uniqueness-for-the-Smooth-Isotropic-Elasticity-Inverse-Problem-September-24-2026/build/sections/introduction.tex}.
\item[S2] OpenAI, manuscript citation page. Source of the formulation motivating this local problem.
\url{https://github.com/openai/math/blob/adc7f1241b42e322a6451854ab7e4b4c146bf78a/preprints/Global-Uniqueness-for-the-Smooth-Isotropic-Elasticity-Inverse-Problem-September-24-2026/README.md}.
\end{itemize}
\end{document}
